\SourceLeaf{321}
oder einer andern (unter den vier jedesmal zusammengehörigen Variationen),
erkennt man dann erst hinterher,\\\hspace*{2em}falls man nicht gleich anfangs mein
Criterium angewandt hat.

Ich setze das Criterium selbst nicht her, um Ihnen das Vergnügen, es selbst
zu finden, nicht zu rauben. Jedenfalls hoffe ich durch das Vorstehende
hinlänglich klar gemacht zu haben, \textls[120]{von was es sich handelt}. Übrigens ist wohl
unnöthig, zu bemerken, dass der ganze Gegenstand höchst elementarisch ist.
Nur ist um so mehr zu verwundern, dass man ihn so ganz übersehen hat.

\mbox{. . . . . . . . . .}

\Needspace{12\baselineskip}\begin{center}[2.]\end{center}

\begin{center}\textsc{Gauss} an \textsc{Gerling}. Göttingen, 24. October 1840.\end{center}

\indent \mbox{. . . . . . . .} Ihre kleine Schrift über die \textsc{Pothenot}sche Aufgabe habe ich
mit Vergnügen gelesen; sie wird allen, die die Aufgabe praktisch anwenden
wollen, sehr lehrreich und nützlich sein. Ich habe mich aber oft gewundert,
dass die vielen Schriftsteller, die sich damit beschäftigt haben, einen in
theoretischer Beziehung ganz wesentlichen Theil ganz übergangen haben,
nemlich die Angabe der mathematischen Bedingungen, unter denen die Aufgabe
\textls[120]{physisch möglich} ist, und die Darstellung dieser Bedingungen in einer
einfachen eleganten Form. Ich mache mich wohl auf folgende Art am
leichtesten verständlich.

Die Rechnungsvorschriften auf vorgegebene numerische Data angewandt geben
(wenn nicht die vier Punkte in Einem Kreise liegen, wo man $\frac{0}{0}$
erhält) immer Ein einziges bestimmtes Resultat. Aber man erkennt leicht,
dass man \textls[120]{einerlei} Resultat erhält, es möge aufgegeben sein, dass der zweite
Punkt von dem ersten um den Winkel $A$ rechts, oder dass er von demselben um
$180^0-A$ links liegt. Eben so beim dritten Punkt. Oder mit andern Worten:

Ist vorgeschrieben, der zweite Punkt solle um $A$ rechts vom ersten liegen,
und der dritte um $B$ rechts vom ersten, so findet man ein Resultat, welches
\textls[120]{dasselbe} ist, was man bei drei andern Datis finden würde, nemlich
